Let $\mathcal{F}$ be the feasible set of atom distributions satisfying the
simplex constraints and all declared evidence:
\[
  \mathcal{F} = \{\pi \in \R^{2^m}: \pi_z \ge 0,\ \sum_z \pi_z = 1,\ A\pi \le b,\ C\pi=d\}.
\]
For a query $\phi$, define
\[
  L_\phi = \inf_{\pi\in\mathcal{F}}\E_\pi[\phi(Z)],
  \qquad
  U_\phi = \sup_{\pi\in\mathcal{F}}\E_\pi[\phi(Z)].
\]

\begin{theorem}[Finite sharpness]
If $\mathcal{F}$ is nonempty, then $L_\phi$ and $U_\phi$ are attained by
feasible atom distributions. Every probability in $[L_\phi,U_\phi]$ is
compatible with the declared linear constraints after mixing endpoint
witnesses.
\end{theorem}

\begin{proof}
The feasible set is a closed and bounded polytope in finite-dimensional
Euclidean space, hence compact. The query functional is linear and therefore
continuous, so it attains its minimum and maximum on $\mathcal{F}$. Convex
mixtures of feasible distributions remain feasible and interpolate the query
value between endpoint optima.
\end{proof}

Classical Frechet-Hoeffding bounds are recovered when the only evidence is
singleton marginals. For an intersection event,
\[
  \Prob(\cap_i \{Z_i=1\}) \in
  \left[\max\left\{0,\sum_i p_i-(m-1)\right\},\ \min_i p_i\right],
\]
and for a union event,
\[
  \Prob(\cup_i \{Z_i=1\}) \in
  \left[\max_i p_i,\ \min\left\{1,\sum_i p_i\right\}\right].
\]
These closed forms are not exponential to compute. Exponential scaling enters
when the general atom-LP path is used for arbitrary Boolean events or side
constraints: the explicit LP has $2^m$ atom variables.
